\documentclass{article} % For LaTeX2e
\usepackage{iclr2026_conference,times}

% Optional math commands from https://github.com/goodfeli/dlbook_notation.
\input{math_commands.tex}

\usepackage{hyperref}
\usepackage{url}

\title{Probing Transformer Reasoning Transfer \\
Through Alphabet Equivalence}

% Double-blind review: authors must not appear in the submitted version.
% \iclrfinalcopy remains commented out until the camera-ready version.
\author{Anonymous authors \\ Paper under double-blind review}

%\iclrfinalcopy % Uncomment for camera-ready version, but NOT for submission.

\begin{document}

\maketitle

\begin{abstract}
A reasoner is \emph{abstract} if it applies a general procedure to a problem's structure, indifferent to how that structure is concretely instantiated; it is \emph{concrete} if what it learned is tied to the particular instantiation trained on. This gives a testable marker: an abstract reasoner should transfer mastery of one problem instance to another. We study equivalence via alphabet renaming, and test transfer using a weight-transplant control -- copying trained per-symbol weights onto the new instance without retraining -- distinguishing genuine reasoning failure from symbols never trained on. We instantiate this with Sudoku, where an equivalent puzzle is any relabeling of its digits. A transformer trained on one such puzzle transfers to an equivalent one no better than an untrained network -- yet the transplant control recovers its original performance exactly. Reasoning, in that narrow sense, is abstract; recognizing that a new instance calls for it is not. We probe this recognition gap three ways. Freezing the shared network and adapting only a new alphabet's parameters fares worse than training an unfrozen model from scratch on the same data: prior exposure to other alphabets does not ease learning a new one. Recombining already-known symbols into unfamiliar combinations yields graded, partial success -- above chance, short of the seamless transfer an abstract reasoner should show, degrading as combinations diverge from training. Replacing the model's fixed, per-alphabet output vocabulary with a mechanism that selects among a puzzle's own visible symbols helps substantially with recombination, but only marginally with transfer to a truly unseen alphabet. These results isolate a specific failure to recognize alphabet-equivalent instances as such, distinct from a general failure to reason, and show several natural interventions narrow, but do not close, this gap.
\end{abstract}

\section{Introduction}

A model that has learned to solve a problem can be said to reason
\emph{abstractly} about it if the procedure it applies operates on the
problem's structure, indifferent to how that structure happens to be
dressed up in a particular presentation. It reasons only
\emph{concretely} if what it has learned is tied to the specific surface
form it was trained on, rather than to the structure underneath. This
distinction is central to any claim that a learned system reasons rather
than merely fits patterns in its training distribution, but such claims
are frequently made without a test that could actually distinguish the
two possibilities. Benchmarks such as SCAN \citep{lake2018generalization}
and COGS \citep{kim2020cogs} test a closely related property --
generalizing to unseen \emph{combinations} of familiar components -- and
have repeatedly found large gaps between in-distribution and
compositional-generalization accuracy in standard neural architectures.
We study a different, narrower notion of generalization: not new
combinations of familiar parts, but the exact same combination presented
under a purely relabeled vocabulary, which isolates a more basic
question -- does the model recognize the relabeled instance as the
problem it already knows how to solve at all.

We propose a concrete, checkable operationalization of this distinction.
An abstract reasoner should treat two different instances of the same
underlying problem interchangeably: having mastered one, it should
transfer to solving the other, since its procedure does not depend on
which particular instantiation it happens to face. We instantiate
``different instances of the same problem'' through the simplest
transform that leaves a problem's structure completely intact while
changing its surface form entirely: renaming its symbol alphabet. Two
problems $P$ and $Q$ are \emph{equivalent} under this definition when
$Q$ is obtained from $P$ purely by relabeling which symbols represent
which underlying roles -- nothing about the structure, constraints, or
solution changes, only the names used to present it.

This gives more than a single empirical test: it gives a general recipe
applicable to any symbolic or combinatorial domain with an analogous
notion of relabeling. The recipe has two parts. First, define
equivalence via a checkable relabeling, and measure whether a model
trained on one instance transfers, unaided, to an equivalent one.
Second -- and this is the part that is easy to get wrong -- use a
targeted control to separate two failure modes that a naive transfer
evaluation cannot distinguish on its own. A model may fail on a new
instance either because its underlying reasoning does not generalize, or
simply because the new instance's symbols were never seen during
training, so any parameters associated with them are untrained noise
regardless of how good the model's reasoning otherwise is. We isolate
the second explanation with a \emph{weight-transplant} control: after
training, we copy the parameters associated with each of $P$'s symbols
directly onto the parameters for $Q$'s corresponding symbols, using the
known-by-construction correspondence, with no further training. This
tells us how the model would perform on $Q$ if it already knew the
correspondence between the two instances' symbols -- precisely the
information a naive zero-shot evaluation cannot provide on its own.

We instantiate this recipe concretely using Sudoku: a well-defined,
verifiable, and cheaply generated combinatorial puzzle whose solution is
invariant to any consistent renaming of its 9 symbols, making it a clean
testbed for this question. Training a transformer to solve Sudoku
puzzles written in one symbol alphabet, we find that it transfers to an
equivalent puzzle in a different alphabet at chance level -- performance
indistinguishable from a network that was never trained at all -- even
though the weight-transplant control recovers its original accuracy
\emph{exactly}. This is the paper's central result: the model's
procedure does reuse across equivalent instances once it is told how
their symbols correspond; what it lacks is any mechanism for discovering
that correspondence on its own. We then probe this recognition gap along
three further axes -- pretraining diversity, an architectural change
that removes a specific structural bottleneck, and training-time data
diversity at a much larger scale -- and show that each narrows, but does
not close, the gap, with a further finding that training diversity
resolves the gap into two distinct sub-problems with different answers:
recombining already-known symbols is solved by sufficient diversity,
while transfer to a genuinely novel symbol is not.

\paragraph{Contributions.}
\begin{itemize}
  \item An operationalized, reusable methodology for testing knowledge
  transfer between equivalent instances of a symbolic problem, including
  a weight-transplant control that isolates genuine reasoning-transfer
  failure from the confound of never-trained parameters.
  \item Applied to Sudoku, we show naive transfer across
  alphabet-equivalent puzzles fails completely -- statistically
  indistinguishable from an untrained network -- while the transplant
  control recovers full performance exactly, demonstrating the failure
  is one of recognition, not of reusable computation (Section~\ref{sec:results}.1).
  \item We show this recognition gap persists across two further natural
  interventions: pretraining diversity across several known alphabets,
  and a redesigned architecture that removes a specific structural
  bottleneck that prevents zero-shot generalization by construction --
  each narrows but does not close the gap (Section~\ref{sec:results}.2--3).
  \item We show training-time data diversity resolves this gap into two
  distinct sub-problems with different answers: recombining
  already-known symbols in novel combinations is essentially solved by
  sufficient training diversity, while transfer to a genuinely novel,
  never-seen symbol remains unsolved by every intervention tested
  (Section~\ref{sec:results}.4).
\end{itemize}

% TODO: roadmap paragraph and exact section cross-references, to be
% finalized once the Results section numbering is settled (pending the
% disentangling experiment's result, see PAPER_PLAN.md Section 5.5).

\section{Related Work}

\paragraph{Neural Sudoku solving.} \citet{palm2018recurrent} introduced
Recurrent Relational Networks, which solve Sudoku via iterative message
passing between cell-level representations on a graph -- the
architectural lineage our encoder-based model, and its optional
iterative-refinement mechanism, belong to. That work, and subsequent
neural combinatorial solvers, focus on solving accuracy within a single
fixed vocabulary; we instead hold solving accuracy on the training
vocabulary fixed near ceiling throughout, and ask what happens under a
symbol relabeling that leaves the underlying combinatorial problem
completely unchanged.

\paragraph{Compositional and systematic generalization.} SCAN
\citep{lake2018generalization} and COGS \citep{kim2020cogs} test whether
sequence models generalize to novel \emph{combinations} of familiar
primitives, and both find persistent, large gaps between in-distribution
and compositional accuracy across a range of architectures. Our question
is narrower and, in a specific sense, easier: we do not ask the model to
combine familiar parts in new ways, only to recognize the exact same
combination it was trained on, wearing a different vocabulary. That we
still find a total failure to recognize this -- indistinguishable from an
untrained network -- suggests the difficulty compositional-generalization
benchmarks document may in part reflect an even more basic failure to
abstract away from surface symbol identity at all.

\paragraph{Symbol binding and variable binding.} \citet{webb2021emergent}
is the closest direct precedent for our framing: their Emergent Symbol
Binding Network uses an indirection mechanism over external memory that
lets learned rules generalize near-perfectly to entirely novel filler
entities, given only a handful of examples, because its computation
operates on relations between representations rather than their
identity. \citet{altabaa2024abstractors} formalize a version of this idea
for Transformers directly, using a relational cross-attention mechanism
architecturally constrained to depend only on similarity between
representations, never their content. Our Section~\ref{sec:results}.3
pointer-based redesign is a much simpler, less architecturally
constrained step in the same direction -- removing a fixed,
identity-keyed output vocabulary rather than building in an explicit
relational bottleneck -- and our results (helping substantially with one
sub-problem, only marginally with the other) suggest the stronger,
more explicit mechanism of \citet{altabaa2024abstractors} may be
necessary where our simpler intervention is not sufficient (see
Section~\ref{sec:results}.4 and the Discussion).

\paragraph{Pointer and copy mechanisms.} Our Section~\ref{sec:results}.3
architecture, which predicts a blank cell's value by selecting among the
symbols visible elsewhere in the same puzzle instance rather than
classifying into a fixed global vocabulary, is a direct application of
Pointer Networks \citep{vinyals2015pointer} and pointer-generator copy
mechanisms \citep{see2017get} to this setting. Both were originally
motivated by output spaces that vary with the input (sorting an input
sequence; naming a rare word from the source document) rather than by
the alphabet-equivalence question we study; we use the same mechanism
for a different purpose, to remove the specific structural obstacle -- a
per-alphabet output vocabulary -- that makes zero-shot transfer to a new
alphabet impossible by construction in the naive classifier design.

\paragraph{Symbolic mechanisms in pretrained language models.}
Contemporaneously with this direction of work, \citet{yang2025emergent}
identify an emergent symbolic architecture inside pretrained large
language models -- symbol abstraction, symbolic induction, and retrieval
heads operating on abstract variables rather than raw tokens -- that
supports abstract reasoning in a related sense. Our study is
complementary: rather than probing an existing large pretrained model,
we train small transformers from scratch under fully controlled data,
isolating the specific alphabet-relabeling manipulation from every other
confound a pretrained model's training history would introduce.

\paragraph{Positioning.} To our knowledge, the specific combination we
use -- defining instance equivalence via a checkable symbol relabeling,
and a weight-transplant control that isolates genuine
reasoning-transfer failure from the confound of never-trained
parameters -- has not been used together in this form, though each piece
draws on the established ideas above.

\section{Problem Setup and Method}
% TODO: not yet drafted. Content exists as prose in REPORT.md's "Common
% setup" section and at publication-appropriate precision in code
% docstrings (sudoku/alphabets.py, sudoku/pointer_dataset.py,
% model/pointer_transformer.py) -- converting this section is mostly
% compression plus a formal-notation pass, not new material. Needs one
% figure (not yet created): a diagram contrasting the two architectures
% (fixed per-alphabet classifier head vs. per-instance pointer
% selection).

\section{Experiments and Results}
\label{sec:results}
% TODO: not yet drafted. Six subsections planned, per PAPER_PLAN.md's
% narrative arc:
%   \ref{sec:results}.1 Naive transfer + transplant control (Exp 0) --
%     complete, written in results/2026-09-07_two_alphabet_sudoku.md
%     and REPORT.md.
%   \ref{sec:results}.2 Diversity alone doesn't close it: few-shot
%     freeze + mixed-alphabet compositional test (Exp 1) -- complete,
%     written in results/2026-09-09_multi_alphabet_experiments.md and
%     REPORT.md.
%   \ref{sec:results}.3 Removing the architectural bottleneck (Exp 2)
%     -- complete, written in REPORT.md.
%   \ref{sec:results}.4 Training-diversity split: compositional vs.
%     genuine zero-shot (Exp 3) -- complete, written in REPORT.md.
%   \ref{sec:results}.5 Disentangling pool size vs. mixing-exposure --
%     ran on Leonardo (job 58017245, training complete: 98.83% val cell
%     accuracy); evaluation jobs 58097892 (zero-shot) and 58097893
%     (mixed-alphabet) submitted, results pending as of this writing.
% Two figures needed (currently zero exist, both buildable now from
% results/*.jsonl without new experiments): (1) per-trial accuracy
% across the mixed-alphabet eval for every condition run so far
% (box/strip plot -- would make the "variance collapse" finding from
% Exp 3 immediate); (2) zero-shot-on-E accuracy/valid\_rate across the
% same conditions (bar or line plot, showing the flat 0\% valid\_rate
% throughout).

\section{Discussion}
% TODO: not yet drafted. Exists as prose in REPORT.md's Synthesis
% section at close-to-paper quality -- needs the disentangling result
% (\ref{sec:results}.5) folded in once available, connection back to
% Section 1's abstract/concrete framing, and a forward-looking paragraph
% on the Abstractors-style relational architecture
% \citep{altabaa2024abstractors} as the natural next step (resolved as
% future work, not a completed experiment in this paper -- see
% PAPER_PLAN.md Open Decisions). Also discusses the (currently
% unverified) vector-equality/duplicate-detection hypothesis for the
% small residual zero-shot signal observed in \ref{sec:results}.3--4,
% explicitly flagged as speculative pending an interpretability probe.

\section{Limitations}
% TODO: not yet drafted. Material exists as scattered caveats throughout
% REPORT.md: single small transformer trained from scratch, single task
% domain (Sudoku-only is a deliberate scope choice, not an oversight --
% see PAPER_PLAN.md Open Decisions); only one equivalence transform
% tested (alphabet relabeling); the pointer architecture's structural
% blind spot on the $\sim$0.8\% of puzzles missing a digit from their
% givens; single seed per major training condition.

\section{Conclusion}
% TODO: not yet drafted; should be straightforward once Discussion is
% finalized.

\bibliography{references}
\bibliographystyle{iclr2026_conference}

\end{document}
